% ============================================================
% 第 16 章 同步电机概述（教材第五篇，印刷页 176--182）
% 节名：第一节 同步电机的基本结构；第二节 同步电机的运行状态；第三节 同步电机的励磁方式；第四节 同步电机的型号和额定值
% ============================================================
\chapter{同步电机概述}

同步电机是电力系统的核心设备：几乎全部交流电能都由同步发电机产生。它的特点是\textbf{转速恒等于同步转速} $n_1=60f/p$，转子转速与旋转磁场"同步"。本章介绍结构与额定值，第 17 章讲基本理论，第 18 章讲并网运行，第 19 章讲电动机与补偿机。

\section{同步电机的基本结构}

同步电机由\textbf{定子}（电枢）、\textbf{转子}（磁极）和气隙组成。定子结构与感应电机基本相同：铁芯由硅钢片叠成，内圆开槽嵌放三相对称电枢绕组。区别主要在转子：同步电机转子装有\textbf{直流励磁的主磁极}。

\subsection{隐极式转子}
转子为圆柱形整体（由整块高强度合金钢锻成），励磁绕组嵌放在转子表面的槽内，槽的分布使气隙磁密接近正弦。特点：
\begin{itemize}
  \item 气隙均匀、转子直径小、轴向长度大；
  \item 机械强度高，可高速旋转（3000\,r/min、1500\,r/min）；
  \item 用于\textbf{汽轮发电机}（火力发电、核电）。
\end{itemize}

\subsection{凸极式转子}
转子磁极是凸出的（磁极铁芯由薄钢板叠成，套上励磁绕组），磁极与磁极之间有较大的极间空隙。特点：
\begin{itemize}
  \item 气隙不均匀（极下小、极间大），$d$、$q$ 轴磁阻不同（这是凸极同步电机分析的出发点）；
  \item 转子直径大、轴向长度短，转速低（几十到几百 r/min）；
  \item 用于\textbf{水轮发电机}及柴油发电机。
\end{itemize}

\begin{center}
\begin{tabular}{@{}p{2.2cm}p{4.2cm}p{4.2cm}@{}}
\toprule
项目 & 隐极式 & 凸极式 \\
\midrule
转子形状 & 圆柱形整体锻钢，槽内嵌励磁绕组 & 凸出磁极，极间空隙大 \\
气隙 & 均匀 & 不均匀（$d$、$q$ 轴磁阻不同） \\
转速 & 高（1500、3000\,r/min） & 低（几十～几百 r/min） \\
极数 & 少（2、4 极） & 多（可达数十极） \\
典型应用 & 汽轮发电机、核电机组 & 水轮发电机、柴油发电机 \\
分析方法 & 隐极机相量图（同步电抗 $X_t$） & 双反应理论（$X_d$、$X_q$） \\
\bottomrule
\end{tabular}
\end{center}

\section{同步电机的运行状态}

同步电机的运行状态由\textbf{功率角}（功角）$\delta$ 的符号与大小决定：$\delta$ 是励磁电动势 $\dot E_0$ 与端电压 $\dot U$ 之间的夹角。

\subsection{发电机状态}
转子由原动机拖动，转速为同步转速 $n_1$，$\dot E_0$ 超前 $\dot U$ 一个正角 $\delta$，电机向电网输出有功功率。这是同步电机最主要的运行状态。

\subsection{电动机状态}
转子拖着负载旋转，从电网吸收有功功率，$\dot E_0$ 滞后 $\dot U$（$\delta<0$），电磁转矩为驱动转矩。同步电动机的转速严格等于 $n_1$，与负载无关（除非失步），故用于要求恒速的场合（如大型压缩机、磨机）。

\subsection{调相机（补偿机）状态}
电机轴上不带机械负载，专门向电网输出或吸收\textbf{无功功率}以改善功率因数。过励时输出无功（相当于电容器），欠励时吸收无功（相当于电抗器）。这是同步电机"第三种用途"。

\begin{center}
\begin{tabular}{@{}llll@{}}
\toprule
运行状态 & 有功流向 & 功角 $\delta$ & 主要用途 \\
\midrule
发电机 & 电机 $\rightarrow$ 电网 & $\delta>0$ & 发电（主要用途） \\
电动机 & 电网 $\rightarrow$ 电机 & $\delta<0$ & 恒速拖动 \\
调相机 & 只交换无功 & $\delta\approx0$ & 改善功率因数、调压 \\
\bottomrule
\end{tabular}
\end{center}

\section{同步电机的励磁方式}

励磁系统为转子提供直流励磁电流，直接决定电机的无功功率与端电压，是同步电机最重要的辅助系统之一。

\begin{enumerate}
  \item \textbf{直流励磁机励磁}：用同轴的直流发电机供电。结构简单、独立性强，但容量受换向限制、维护量大，已逐渐淘汰。
  \item \textbf{静止半导体励磁}：由励磁变压器从机端取电，经可控整流后供给转子。响应快、维护简单，是大中型机组的主流。
  \item \textbf{旋转半导体（无刷）励磁}：交流励磁机与主发电机同轴，其电枢（旋转）输出的交流经装在转轴上的旋转整流器整流后直接供主发电机励磁，无需滑环电刷。可靠性高，用于大型机组。
  \item \textbf{三次谐波励磁}：利用气隙磁场中的三次谐波在附加绕组中感应电动势，经整流后励磁，具有自励恒压特性，多用于小型机组。
\end{enumerate}

\section{同步电机的型号和额定值}

\subsection{额定容量 $S_N$ 与额定功率 $P_N$}
同步发电机铭牌上给出\textbf{额定视在功率} $S_N$（kV·A）和\textbf{额定有功功率} $P_N$（kW），二者关系为
\begin{equation}
P_N=S_N\cos\varphi_N .
\end{equation}
对于同步电动机，额定功率指轴上输出的机械功率；对于调相机，只给额定视在功率。

\subsection{额定电压 $U_N$}
额定运行时定子绕组的\textbf{线电压}，单位 V 或 kV。

\subsection{额定电流 $I_N$}
额定运行时定子绕组的\textbf{线电流}，单位 A。由
\begin{equation}
S_N=\sqrt3U_NI_N
\end{equation}
确定。

\subsection{额定功率因数 $\cos\varphi_N$}
额定运行时的功率因数，一般发电机为 0.8（滞后），电动机为 0.8～0.9（超前或滞后按铭牌）。

\subsection{额定频率 $f_N$}
我国为 50\,Hz。极对数 $p$ 与转速严格对应：$n_N=60f_N/p$，故同步电机转速是"分档"的（3000、1500、1000、750、600…r/min）。

\subsection{额定励磁电流 $I_{fN}$ 与额定励磁电压 $U_{fN}$}
额定运行（额定电压、额定电流、额定功率因数）时的转子励磁电流与电压。

\keyformula{S_N=\sqrt3U_NI_N,\qquad P_N=S_N\cos\varphi_N,\qquad n_N=\frac{60f_N}{p}}

\corepoint{同步电机额定值有两处与感应电机不同：① 发电机铭牌给的是\textbf{视在功率} $S_N$（因为要同时保证电压和功率因数），电动机才给轴上机械功率 $P_N$；② 转速不是"额定"出来的，而是由极对数和频率\textbf{唯一确定}（$n=f/p$ 的整数分档），所以同步电机没有"额定转差率"这一说。}
